% ============================================================
\part{Git und GitHub}
% ============================================================

\chapter{Wir wollen einen Fehler rückgängig machen}

\kapitelziel{
Du verstehst die Git-History von robotik-uno-q, kannst
neue Commits sinnvoll schreiben und Änderungen rückgängig machen.
}
\kapitelstruktur

\section{Situation}

\texttt{robotik-uno-q} ist bereits ein Git-Repository --
du hast es von GitHub geklont. Es hat eine History.
Du hast seitdem Dateien hinzugefügt: \texttt{sensor\_bridge.py},
\texttt{sensor\_config.yaml}, \texttt{start\_roboter.sh}.

Diese Änderungen sind noch nicht in Git.
Wenn jetzt etwas schiefgeht -- die Festplatte stirbt,
ein Update zerstört den Code -- ist alles weg.

\begin{aufgabeBox}
Versioniere alle bisherigen Änderungen am Projekt sinnvoll,
sodass du bei Bedarf auf jeden Stand zurückspringen kannst.
\end{aufgabeBox}

\begin{haubeBox}
\textbf{Warum so viele kleine Commits statt einem großen?}

Stell dir vor, du hast drei Wochen Arbeit in einen einzigen
Commit gepackt: \enquote{Projekt fertig.} Jetzt funktioniert
etwas nicht mehr -- aber was? Du musst mühsam im Gedächtnis
kramen, was du wann geändert hast.

Mit kleinen Commits ist jede Änderung einzeln rückgängig
zu machen. \texttt{git revert abc1234} macht genau diesen
einen Commit rückgängig -- alle anderen bleiben.
Das ist wie \enquote{Strg+Z} für einzelne Entscheidungen.
\end{haubeBox}

\section{Der pragmatische Weg}

\begin{lstlisting}[style=bash]
cd ~/robotik-uno-q

# Was hat sich geaendert?
git status
git diff

# Aenderungen versionieren
git add stufe2_ros2/src/sensor_bridge.py
git commit -m "feat(stufe2): sensor_bridge.py -- Grundstruktur"

git add stufe2_ros2/config/sensor_config.yaml
git commit -m "feat(config): HC-SR04 Sensor-Konfiguration"

git add scripts/start_roboter.sh
git commit -m "feat(scripts): start_roboter.sh -- Stack-Startskript"

# History anzeigen
git log --oneline
\end{lstlisting}

\begin{loesungBox}
\textbf{Fehler rückgängig machen:}
\begin{lstlisting}[style=bash]
# Letzte Aenderungen in sensor_bridge.py verwerfen (vor Commit)
git checkout -- stufe2_ros2/src/sensor_bridge.py

# Letzten Commit rueckgaengig (neuen Commit erstellt)
git revert HEAD

# Auf bestimmten Stand zurueckgehen (nur lesen, nicht zerstoeren)
git show abc1234:stufe2_ros2/src/sensor_bridge.py > /tmp/alte_version.py
\end{lstlisting}
\end{loesungBox}

\begin{projektBox}
\textbf{Bisher:} Ungekoppelte Dateien.\\
\textbf{Heute:} Saubere Git-History mit aussagekräftigen Commit-Nachrichten.\\
\textbf{Nächstes Kapitel:} Sicher experimentieren mit Branches.
\end{projektBox}

\begin{merksatzBox}
\textbf{Commit-Nachrichten-Format:} \texttt{typ(bereich): was wurde gemacht}

\texttt{feat}: neues Feature ·
\texttt{fix}: Bugfix ·
\texttt{docs}: Dokumentation ·
\texttt{refactor}: Umbau ohne neue Funktion

\enquote{Fix} ist eine schlechte Nachricht.
\enquote{fix(serial): Timeout bei langsamer Arduino-Antwort erhöht} ist eine gute.
\end{merksatzBox}

\begin{robotikBox}
ROS2-Pakete auf GitHub folgen Conventional Commits.
Wer sich jetzt angewöhnt sauber zu committen, kann
später Pull Requests für offizielle ROS2-Pakete einreichen --
das Format ist identisch.
\end{robotikBox}

% ── Kapitel 14 ────────────────────────────────────────────────
\chapter{Wir wollen gefahrlos etwas ausprobieren}

\kapitelziel{
Du kannst Branches für neue Features nutzen und damit
\texttt{stufe1\_mcu/} und \texttt{stufe2\_ros2/} parallel entwickeln.
}
\kapitelstruktur

\section{Situation}

\texttt{hinderniserkennung.ino} funktioniert. Jetzt willst
du eine verbesserte Filterung ausprobieren -- Mittelwert
über die letzten 5 Messungen statt Direktwert.
Aber was, wenn die neue Version schlechter ist?

\begin{aufgabeBox}
Erstelle einen Branch für die verbesserte Hinderniserkennung
und entwickle dort, ohne die funktionierende Version zu gefährden.
\end{aufgabeBox}

\begin{lstlisting}[style=bash]
# Branch erstellen
git checkout -b feature/ultraschall-mittelwertfilter

# Aenderung vornehmen
nano stufe1_mcu/hinderniserkennung/hinderniserkennung.ino
# Mittelwert-Filterung einbauen

git add stufe1_mcu/hinderniserkennung/hinderniserkennung.ino
git commit -m "feat(stufe1): Ultraschall -- gleitender Mittelwert (n=5)"

# Testen -- funktioniert es besser?
# Wenn ja: zurueck zu main und mergen
git checkout main
git merge feature/ultraschall-mittelwertfilter

# Wenn nein: Branch einfach liegen lassen oder loeschen
git branch -d feature/ultraschall-mittelwertfilter
\end{lstlisting}

\begin{merksatzBox}
\texttt{main} ist immer stabil. Neue Ideen kommen auf
Feature-Branches. Ein Merge passiert erst, wenn die
neue Funktion getestet und überlegen ist.
\end{merksatzBox}

% ── Kapitel 15 ────────────────────────────────────────────────
\chapter{Wir wollen unser Projekt sichern und teilen}

\kapitelziel{
Du kannst deine Änderungen zu GitHub pushen und das Repo
mit anderen teilen -- auch mit dem Raspberry Pi des Roboters.
}
\kapitelstruktur

\section{Situation}

Alle Änderungen sind lokal committed. Wenn der Laptop
abstürzt, ist alles weg. Und der Raspberry Pi des Roboters
hat noch keinen der neuen Dateien.

\begin{aufgabeBox}
Push die neuen Commits zu \texttt{wollbaer1971/robotik-uno-q}
und richte auf dem Raspberry Pi einen automatischen Pull ein.
\end{aufgabeBox}

\begin{lstlisting}[style=bash]
# Aenderungen hochladen
git push origin main

# Auf dem Raspberry Pi (per SSH):
ssh ros@192.168.1.100
cd ~/robotik-uno-q
git pull
# stufe2_ros2/src/sensor_bridge.py ist jetzt auch auf dem Pi
\end{lstlisting}

\begin{lstlisting}[style=bash]
# Noch besser: Automatischer Pull beim Start
# In start_roboter.sh ganz oben:
echo "Aktualisiere Projektdateien..."
cd $PROJEKT
git pull --quiet
\end{lstlisting}

\begin{projektBox}
\textbf{Bisher:} Lokale Git-History.\\
\textbf{Heute:} Alle Änderungen auf GitHub -- Laptop-Absturz
kostet nur eine Stunde, nicht drei Wochen.\\
\textbf{Raspberry Pi:} Bei jedem Start automatisch aktuell.\\
\textbf{Nächstes Kapitel:} Docker für reproduzierbare Umgebungen.
\end{projektBox}

\begin{merksatzBox}
\textbf{Täglicher Workflow:}
1. \texttt{git pull} -- immer zuerst.
2. Arbeiten, kleine Commits schreiben.
3. \texttt{git push} -- am Ende des Tages oder nach jedem Feature.
GitHub ist das Backup. Der Raspberry Pi zieht automatisch die
neueste Version bei jedem Start.
\end{merksatzBox}

\begin{robotikBox}
Professionelle Robotikprojekte haben oft einen
\enquote{Release}-Branch, der auf den Robotern läuft,
und einen \enquote{develop}-Branch für laufende Entwicklung.
Der Roboter führt automatisch \texttt{git pull} auf
\texttt{release} aus -- niemals auf \texttt{develop}.
Das verhindert, dass ein halbfertiges Feature den
laufenden Roboter zum Absturz bringt.
\end{robotikBox}
